\section*{Chapter \ref{ch:pl:accelerators-for-quantitative-work} --- Accelerators for Quantitative Work}

\begin{solution}{exo:pl:accelerators-for-quantitative-work:1}
$2n^3$ operations over $3 \times 8n^2$ bytes: $n/12$. It is \omterm{def:pl:accelerators-for-quantitative-work:bound}{compute-bound} when $n/12$ exceeds the ridge: from $n = 34$ on the core (ridge 2.79),
from $n = 122$ on the device (ridge 10.1).
\end{solution}

\begin{solution}{exo:pl:accelerators-for-quantitative-work:2}
Its intensity (below 1.3 operations per byte for every kernel) is below both ridges, so on both machines its time is bytes over bandwidth,
and the ratio of times is the ratio of bandwidths, 140; the device's much larger peak is never reached.
\end{solution}

\begin{solution}{exo:pl:accelerators-for-quantitative-work:3}
$1/(0.05 + 0.95/140) = 17.6$; with 99\%, $1/(0.01 + 0.99/140) = 58.6$. The part that stays on the host sets the ceiling: 20 and 100 at
infinite speed-up.
\end{solution}

\begin{solution}{exo:pl:accelerators-for-quantitative-work:4}
The roofline counts \texttt{exp} as one operation, but an exponential takes many cycles; the kernel is limited by the rate of the
exponential, not by bandwidth. A vectorised exponential (a library with SIMD transcendental functions), computing the exponential of the
cumulative sum once per path instead of per step, or fusing the step with the next kernel so that $S$ is not written and read again, would
move it.
\end{solution}

\begin{solution}{exo:pl:accelerators-for-quantitative-work:5}
Every trade's kilobyte crosses the link at 128~GB/s (7.8~ns) and its 8~kB of memory traffic costs the device 2.4~ns; the CPU needs 334~ns per
trade. The ratio of the CPU's time per trade to the device's is 32.7, whatever the batch, once the launch latency is amortised.
\end{solution}

\begin{solution}{exo:pl:accelerators-for-quantitative-work:6}
The path generation and the payoff at each date run the same instructions on every path: no divergence. The exercise decision (continue or
exercise, depending on the path's value against the regression) branches by path, and paths already exercised stop contributing; the
regression itself is a set of reductions, fine on the device but with a synchronisation at every date.
\end{solution}

\begin{solution}{exo:pl:accelerators-for-quantitative-work:7}
A megabyte per request adds 7.8~microseconds of transfer to the 10 of latency: the break-even rises from 31 to 56 trades per request. Keeping
market data resident on the device and sending only what changed is worth as much as halving the launch latency.
\end{solution}

\begin{solution}{exo:pl:accelerators-for-quantitative-work:8}
The peak ratio applies only to \omterm{def:pl:accelerators-for-quantitative-work:bound}{compute-bound} work; a risk run of Monte Carlo and aggregation is \omterm{def:pl:accelerators-for-quantitative-work:bound}{memory-bound} and gains at most the bandwidth
ratio, 140 (four minutes, not one), before transfers, launches and the host part (Amdahl). And the comparison is with one core: against the
server's whole socket the gain is several times smaller.
\end{solution}

\begin{solution}{pb:pl:accelerators-for-quantitative-work:1}
\begin{enumerate}
\item 66.9~GFLOP/s, 24.0~GB/s, a ridge of 2.79 operations per byte.
\item 34~TFLOPS in double precision, 3.35~TB/s, a ridge of 10.1; 508 and 140 times the core.
\item The path step, the basket payoff and the exposure aggregation on both; the matrix product is \omterm{def:pl:accelerators-for-quantitative-work:bound}{compute-bound} on both.
\item It is limited by the exponential, which the counting convention treats as one operation.
\item Latency, transfers, divergence, precision, porting cost, price and power.
\item 66 billion operations and 98~GB of memory traffic; only the trades' descriptions go over the link and the profile comes back.
\item 4.09~seconds on the core, 29.3~milliseconds on the device.
\item It is \omterm{def:pl:accelerators-for-quantitative-work:bound}{memory-bound}: 140 is the bandwidth ratio.
\item At most 17.6 times if 95\% of the process moves.
\item The per-core ratio falls to the ratio of the device's bandwidth to the socket's.
\item 0.33~microseconds on the core, 10.0 on the device: 30 times slower.
\item At 31 trades per request.
\item 4, 16, 31 and 155 trades for launch latencies of 1, 5, 10 and 50~microseconds.
\item The link: 32.7 times at most.
\item Batch its requests (hundreds per call) and keep market data on the device.
\item \textbf{Named result.} The exposure Monte Carlo runs 140 times faster end to end on the device than on one core (17.6 times for the
whole process at 95\% offload); a single-trade price request runs 30 times slower; the device breaks even at 31 trades per request with a
10-microsecond launch latency.
\item Large, batched, memory- or \omterm{def:pl:accelerators-for-quantitative-work:bound}{compute-bound} array work with little data crossing: exposure and XVA Monte Carlo, scenario revaluation,
model training.
\item Its sustained bandwidth and double-precision rate on our kernels, the launch and transfer latency of our requests, and the porting
effort of one real kernel.
\item Because risk and pricing results are compared and reconciled at double precision; single precision would change the device's roof and
must be validated first.
\item When the work is large, batched and stays on the device long enough to amortise what crossing to it costs.
\end{enumerate}
\end{solution}

\begin{solution}{iq:pl:accelerators-for-quantitative-work:1}
Operations per byte moved. Below a machine's ridge point a kernel is limited by bandwidth, above it by arithmetic; hardware with more peak
does nothing for a \omterm{def:pl:accelerators-for-quantitative-work:bound}{memory-bound kernel}, and the ratio of bandwidths is the most it can gain from a faster memory.
\iqlookfor{Roofline reasoning.}
\end{solution}

\begin{solution}{iq:pl:accelerators-for-quantitative-work:2}
The Monte Carlo is large and amortises the device's fixed costs; each pricing request is small, so its time is dominated by the launch,
synchronisation and transfer, which the CPU does not pay. Batch requests or keep them on the CPU.
\iqlookfor{Fixed costs against work size.}
\end{solution}

\begin{solution}{iq:pl:accelerators-for-quantitative-work:3}
Count the kernel's operations, memory bytes and transferred bytes; take the device's peak, bandwidth and link from its datasheet and a
launch latency from a measurement or a conservative assumption; bound its time with the roofline plus fixed costs; compare with the CPU's
measured roof; apply Amdahl's law to the whole process.
\iqlookfor{A cost model with stated assumptions.}
\end{solution}

\begin{solution}{iq:pl:accelerators-for-quantitative-work:4}
With a counter-based generator the random numbers depend only on the path, step and seed, so both machines draw the same ones; the
arithmetic must then be done in the same precision and order (reductions are the usual difference), and transcendental functions compared to a
tolerance or taken from the same library.
\iqlookfor{Counter-based generators; reductions and precision.}
\end{solution}

\begin{solution}{iq:pl:accelerators-for-quantitative-work:5}
Inventory the workloads with their operations, bytes, transfers and batch sizes; place them on measured and cited rooflines; predict end-to-end
gains with the fixed costs and Amdahl's law; port one kernel of the most promising workload and measure; decide on the measured gain against
porting, maintenance, price and power, and keep a CPU implementation as the reference.
\iqlookfor{Model first, measure one kernel, keep the reference.}
\end{solution}
